\documentclass[12pt,a4paper]{article}
\usepackage{amsmath}
\usepackage{amssymb}
\usepackage{amsthm}
\usepackage{geometry}
\usepackage{xcolor}
\usepackage{hyperref}

\geometry{margin=1in}

% Define theorem environments
\theoremstyle{definition}
\newtheorem{definition}{Definition}[section]
\newtheorem{theorem}{Theorem}[section]
\newtheorem{proposition}{Proposition}[section]
\newtheorem{lemma}{Lemma}[section]
\newtheorem{remark}{Remark}[section]
\newtheorem{example}{Example}[section]

\title{COMP0086: Bayesian Deep Learning - Lecture Notes}
\author{Comprehensive Course Notes}
\date{}

\begin{document}

\maketitle
\tableofcontents
\newpage

\section*{Introduction}

This course covers the intersection of Bayesian inference and deep learning, exploring uncertainty quantification, efficient computation, and the connection between probabilistic inference and neural network learning.

\newpage

\section{Part I: Foundational Optimization and Learning}

\subsection{Regularization and Bayesian Priors}

A regularizer $r(w)$ of order $p$ is:
\[
r(w) = \sum_i |w_i|^p, \quad p > 0
\]

The $L^2$ regularization (Ridge) corresponds to a Gaussian prior. The $L^1$ regularization (Lasso) corresponds to a Laplace prior. The Laplace distribution is:
\[
p(w) = \frac{1}{2b} \exp\left(-\frac{|w-\mu|}{b}\right)
\]

\subsection{Loss Functions for Classification}

For binary classification:
\[
\sum_i y_i \log f(x_i, w) + (1-y_i) \log(1 - f(x_i, w))
\]

For multi-class classification with softmax:
\[
\hat{y}_{ik} = \text{softmax}(z_i) = \frac{\exp(z_{ik})}{\sum_j \exp(z_{ij})}
\]

The loss function is:
\[
\ell(z_i, y_i) = -\sum_{j=1}^K y_{ij} \log \hat{y}_{ij}
\]

\subsection{Gradient Descent and Stochastic Optimization}

The full loss is:
\[
L(\theta) = \frac{1}{n} \sum_{i=1}^n L_i(\theta) = \frac{1}{n} \sum_{i=1}^n \ell(f(x_i, \theta), y_i)
\]

Stochastic gradient descent estimates using mini-batches:
\[
\nabla_\theta L(\theta) \approx \frac{1}{m} \sum_{j=1}^m \nabla_\theta \ell(f(x_j[\theta], y_j))
\]

Small batches lead to faster updates; large batches provide better gradient estimates.

\vspace{0.3cm}
\noindent\textbf{Moving Average:} The moving average at step $t$ is updated as:
\[
\bar{a}_t = \frac{1}{t} \sum_{k=1}^t x_k = \epsilon_t x_t + \mu_t \bar{a}_{t-1}
\]

where if $\mu_t$ is small, more recent samples contribute more strongly.

\subsection{Momentum and Adaptive Methods}

\textbf{Momentum-based updates} accelerate convergence by accumulating gradients:
\[
g_{t+1} = \mu g_t - \nabla f(\theta_t), \quad \theta_{t+1} = \theta_t + g_{t+1}
\]

\textbf{AdaGrad} adapts learning rates per parameter using accumulated squared gradients:
\[
\nabla_i = \nabla f(\theta)_i, \quad G_{t, ii}^t = \sum_{k=1}^t (\nabla_i)^2, \quad \theta_t^m = \theta_{t-1}^m - \frac{\eta}{\sqrt{G_{t,tt}^m + \epsilon}} g_t^m
\]

However, accumulated gradients only increase, causing learning rates to decay to zero.

\textbf{RMSprop} uses a moving average of squared gradients to solve this:
\[
S_{t,i} = \gamma S_{t-1,i} + (1-\gamma) g_{t,i}^2, \quad \theta_t^m = \theta_t^m - \frac{\sqrt{S_{t-1,i} + \epsilon}}{g_{t,i}}
\]

\textbf{Adam Optimizer} combines RMSprop with Momentum and adds bias correction:
\[
m_{t,i} = \beta_1 m_{t-1,i} + (1-\beta_1) g_{t,i}, \quad \hat{m}_{t,i} = \frac{m_{t,i}}{1-\beta_1^t}
\]
\[
s_{t,i} = \beta_2 s_{t-1,i} + (1-\beta_2) g_{t,i}^2, \quad \hat{s}_{t,i} = \frac{s_{t,i}}{1-\beta_2^t}
\]
\[
\theta_t^m = \theta_t^m - \frac{\eta}{\sqrt{\hat{s}_{t,i}} + \epsilon} \hat{m}_{t,i}
\]

\subsection{Quasi-Newton Methods}

Higher-order optimization methods find $B$ satisfying the secant condition:
\[
\nabla f(\theta_{t+1}) = \nabla f(\theta_c) + B(\theta_{t+1} - \theta_c)
\]

These achieve superlinear convergence order.

\newpage

\section{Part II: Automatic Differentiation and Computation}

\subsection{Computing Gradients Efficiently}

Automatic differentiation computes gradients via the chain rule. There are two principal modes:

\begin{itemize}
  \item \textbf{Forward mode:} Computes Jacobian-vector products. Memory-efficient to implement but best when the number of inputs is much smaller than outputs.
  \item \textbf{Backward mode:} Computes vector-Jacobian products. More memory-intensive but efficient when inputs $\gg$ outputs (common in deep learning).
\end{itemize}

\textbf{Forward propagation:} $\dot{u}_{t+1} = \frac{\partial u_{t+1}}{\partial u_t} \dot{u}_t$

\textbf{Backward propagation:} $\bar{V}_n = D_n \bar{V}_{n+1}$

\subsection{Dual Numbers and Forward Mode}

Dual numbers extend the reals with an infinitesimal $\epsilon$ where $\epsilon^2 = 0$:
\[
f(x + \epsilon) = f(x) + f'(x) \epsilon
\]

Forward-mode automatic differentiation uses dual number arithmetic to compute derivatives while evaluating the function.

\subsection{Handling Constraints via Transformations}

For constrained parameters (e.g., positive variances), use unconstrained representations:
\[
\sigma = \exp(s), \quad s \in \mathbb{R}
\]

For structured covariance matrices, use Cholesky or similar factorizations:
\[
\Sigma = A A^T + \text{diag}(e^s)
\]

This reduces parameters from $O(D^2)$ to $O(Dr + D)$ for rank $r$.

\subsection{Minibatching with Priors}

When using minibatches, scale the prior term to match the dataset size:
\[
L(w) = \log p(w) + \sum_{i=1}^n \log p(z_i | w, x_i)
\]

For minibatch $M$ of size $N/M$:
\[
L_M(w) = \log p(w) + \frac{N}{M} \sum_{i=1}^M \log p(y_i | w, x_i)
\]

This ensures unbiased gradient estimates for the full objective.

\newpage

\section{Part III: Bayesian Inference and Posterior Approximation}

\subsection{Posterior Approximation with Laplace}

The Laplace approximation approximates the posterior with a Gaussian distribution centered at the MAP estimate:
\[
\log p(w|D) \approx \log \mathcal{N}(w | w^*, \Lambda^{-1}) = -\frac{1}{2}(w - w^*)^T \Lambda (w - w^*) + \text{const}
\]

where $w^* = w_{\text{map}}$ and $\Lambda = -\nabla^2 \log p(w|D)$ is the Hessian at the MAP.

The marginal likelihood with Laplace approximation is:
\[
p(D) \approx p(D|\theta^*) p(\theta^*) \frac{(2\pi)^{D/2}}{|\Lambda|^{1/2}}
\]

Broader posteriors (larger determinant) are preferred for model selection.

\subsection{Variational Inference and ELBO}

The evidence lower bound (ELBO) decomposes the log marginal likelihood:
\[
\mathcal{L}(w) = \mathbb{E}_{q(w)} \left[\log \frac{p(w,D)}{q(w)}\right] = \log P(D) - \mathrm{KL}[q(w) \| p(w|D)]
\]

Maximizing the ELBO minimizes the KL divergence to the true posterior. The gradient with respect to variational parameters $\lambda$ is:
\[
\nabla_\lambda \mathbb{E}_{q(w)} \left[\log p(w, D) - \log q_\lambda(w)\right] = \mathbb{E}_{q(w)} \left[\nabla_\lambda \log \frac{p(w,D)}{q_\lambda(w)}\right]
\]

\subsection{Reparametrization Trick}

The reparametrization trick enables gradient estimation through sampling. With $\epsilon \sim \mathcal{N}(0, I)$ and $w = \mu + \Sigma^{1/2} \epsilon$:

\[
\nabla_\lambda \mathbb{E}_{q(w)} \left[\log \frac{p(w,D)}{q_\lambda(w)}\right] = \mathbb{E}_{p(\epsilon)} \left[\nabla_\lambda \log \frac{p(\mu + \Sigma^{1/2}\epsilon, D)}{q_\lambda(\mu + \Sigma^{1/2}\epsilon)}\right]
\]

This decouples the gradient from the expectation, allowing backpropagation through sampling.

\subsection{Understanding KL Divergence Modes}

\textbf{Reverse KL (Mode-seeking):} $\mathrm{KL}[q \| P]$ penalizes placing probability mass where $p$ is zero:
\[
\mathbb{E}_{q(z)} \left[\log \frac{q(z)}{p(z|D)}\right] = \mathbb{E}_{q(z)} \log q(z) - \mathbb{E}_{q(z)} \log p(z|D)
\]

When $p(z|D) \to 0$ but $q(z) > 0$, the cross-entropy diverges to $-\infty$, heavily suppressing wide tails.

\textbf{Forward KL (Mass-covering):} $\mathrm{KL}[p \| q]$ penalizes missing probability mass:
\[
\mathbb{E}_{p(z|D)} \left[\log \frac{p(z|D)}{q(z)}\right] = \mathbb{E}_{p(z|D)} \log p(z|D) - \mathbb{E}_{p(z|D)} \log q(z)
\]

When $p(z|D) > 0$ but $q(z) \approx 0$, it diverges to $\infty$, forcing broad coverage.

\subsection{Sequential Learning}

In streaming or continual learning settings, the posterior from the previous batch becomes the prior for the next:
\[
p(\theta | D_{1:t}) \propto p(D_t | \theta) p(\theta | D_{1:t-1})
\]

This enables efficient online updates without reprocessing historical data.

\newpage

\section{Part IV: Deep Learning Architectures}

\subsection{Deep Feature Learning}

Rather than relying on hand-crafted features, deep learning learns feature representations:
\[
f(x) = w^T \tilde{\Phi}(x)
\]

where both the weights $w$ and the feature transformation $\tilde{\Phi}$ are learned jointly. The simplest deep model has a single hidden layer:
\[
f(x) = w^T g(Ax + b)
\]

with non-linear activation $g(\cdot)$.

\subsection{Multi-layer Perceptrons}

A multi-layer network stacks multiple transformations:
\[
h_0 = x, \quad h_\ell = g(U_\ell h_{\ell-1} + b_\ell), \quad \ell = 1, \ldots, L
\]

where $g(\cdot)$ is typically ReLU: $g(z) = \max(0, z)$.

Loss functions align with probabilistic models:
\begin{itemize}
  \item Squared error $\Leftrightarrow$ Gaussian likelihood
  \item Binary cross-entropy $\Leftrightarrow$ Bernoulli likelihood  
  \item Absolute error $\Leftrightarrow$ Laplace likelihood
\end{itemize}

Regularization via $L^1$ or $L^2$ penalties with early stopping prevents overfitting.

\subsection{Convolutional Neural Networks}

Convolutions exploit two key properties of natural signals:
\begin{enumerate}
  \item \textbf{Locality:} Features depend only on small neighborhoods
  \item \textbf{Translation Invariance:} Patterns repeat across space
\end{enumerate}

The 1D convolution is:
\[
(f * g)(t) = \int f(a) g(t - a) da, \quad \text{or discretely:} \quad (X * w)_i = \sum_{j=0}^{k-1} X_{i+j} w_j
\]

The 2D case extends naturally:
\[
(X * K)_{ij} = \sum_a \sum_b X_{i+a,j+b} K_{ab}
\]

A typical convolutional block consists of:
\begin{itemize}
  \item Convolutional stage: Affine transform with weight sharing
  \item Detector stage: Nonlinearity (e.g., ReLU)
  \item Pooling stage: Max or average pooling over spatial regions
\end{itemize}

\subsection{Network Depth and Architecture Improvements}

\textbf{Residual Connections} enable very deep networks by using skip connections:
\[
h_\ell = h_{\ell-1} + M(h_{\ell-1})
\]

This prevents vanishing gradients and enables information to flow across many layers.

\textbf{Batch Normalization} rescales activations to zero mean, unit variance per batch:
\[
\mu_d = \frac{1}{M} \sum_i x_i^d, \quad \hat{x}_i^d = \frac{x_i^d - \mu_d}{\sqrt{\sigma_d^2 + \epsilon}}, \quad y_i^d = \gamma_d \hat{x}_i^d + \beta_d
\]

Learnable scale and shift parameters $\gamma_d, \beta_d$ restore expressiveness. BN reduces internal covariate shift and enables faster training.

\subsection{Bayesian Approaches to Deep Models}

\textbf{Bayesian Last Layer:} A practical middle ground between point estimates and full Bayesian inference:
\[
p(y|x) = \int p(y | w^T \phi(x)) p(w|D) dw
\]

Find point estimates for early layers and infer posterior over the final layer weights.

\textbf{Deep Ensembles:} Train $S$ networks with different initializations and treat as an ensemble:
\[
p(y|x) \approx \frac{1}{S} \sum_{s=1}^S p(y|x, \theta^{(s)})
\]

Naturally captures multimodal posterior structures and weight permutation symmetries that challenge single-model Bayesian methods.

\textbf{Bayesian Inference Challenges:} Neural networks exhibit many symmetries (e.g., permutation of hidden units), creating disconnected modes in the posterior. Laplace approximation typically captures only one mode, while variational inference is often mode-seeking, and MCMC may mix poorly. Ensembles address this naturally.

\newpage

\section{Part V: Uncertainty Quantification and Probabilistic Models}

\subsection{Uncertainty in Neural Networks}

Instead of point estimates $\theta^*$, we can assign distributions $p(\theta|D)$ to capture parameter uncertainty. The posterior predictive distribution is:
\[
p(y|x, D) = \int p(y|x, \theta) p(\theta|D) d\theta
\]

For Gaussian likelihoods, this decomposes into epistemic (parameter) and aleatoric (data) uncertainty.

\subsection{Laplace Approximation for Deep Networks}

Find the MAP estimate: $\theta^* = \operatorname*{argmax}_\theta \log p(\theta|D)$

Approximate the posterior with a Gaussian:
\[
p(\theta|D) \approx \mathcal{N}(\theta | \theta^*, \mathbb{H}^{-1})
\]

where the Hessian has the form:
\[
\mathbb{H} = \alpha I + \beta \sum_i [\nabla_\theta f(x_i, \theta)] [\nabla_\theta f(x_i, \theta)]^T
\]

Predictive uncertainty combines epistemic and aleatoric components:
\[
p(y|x,D) \approx \mathcal{N}(y | f_{\theta^*}(x), [\nabla f_{\theta^*}(x)]^T \mathbb{H}^{-1} [\nabla f_{\theta^*}(x)] + \sigma^2)
\]

\subsection{Model Calibration}

A classifier is well-calibrated when predicted probabilities match empirical frequencies:
\[
P(\text{true label} = k \mid \text{predicted prob} = p) \approx p
\]

For evaluation, binning predictions into intervals $I_m = [c_{m-1}, c_m]$:
\begin{itemize}
    \item Accuracy: $\text{acc}_m = \frac{1}{|B_m|} \sum_{i \in B_m} \mathbb{1}[\hat{y}_i = y_i]$
    \item Confidence: $\text{conf}_m = \frac{1}{|B_m|} \sum_{i \in B_m} \hat{p}_i$
\end{itemize}

Expected Calibration Error (ECE) measures the gap:
\[
\text{ECE} = \sum_m \frac{|B_m|}{N} |\text{acc}_m - \text{conf}_m|
\]

\textbf{Temperature Scaling} adjusts predicted probabilities without changing decisions:
\[
p_{\text{scaled}}(y=k|x) = \frac{\exp(z_k / T)}{\sum_j \exp(z_j / T)}
\]

Find optimal $T$ by maximizing likelihood on a validation set.

\subsection{Epistemic vs. Aleatoric Uncertainty}

Total predictive uncertainty decomposes into:
\[
p(y | x, D) = \int p(y | x, \theta) p(\theta | D) d\theta
\]

\begin{itemize}
    \item \textbf{Epistemic uncertainty} (model/parameter uncertainty): Reduces with more data. Captures $\mathrm{Var}_\theta[\mathbb{E}[y|\theta]]$
    \item \textbf{Aleatoric uncertainty} (data/measurement uncertainty): Irreducible. Captures $\mathbb{E}_\theta[\mathrm{Var}(y|\theta)]$
\end{itemize}

For Gaussian models:
\[
\mathrm{Var}(y) = \mathrm{Var}_\theta[\mu(\theta)] + \mathbb{E}_\theta[\sigma^2(\theta)]
\]

\subsection{Variational Autoencoders}

VAEs define a generative model through latent variables $z$:
\[
p_\theta(x) = \int p_\theta(x|z) p(z) dz, \quad p(z) = \mathcal{N}(0, I)
\]

The ELBO decomposes into reconstructive and regularization terms:
\[
\log p_\theta(x) \geq \underbrace{\mathbb{E}_{q_\phi(z|x)}[\log p_\theta(x|z)]}_{\text{Reconstruction}} - \underbrace{D_{\mathrm{KL}}(q_\phi(z|x) \| p(z))}_{\text{Regularization}}
\]

The encoder and decoder are typically parameterized as:
\[
q_\phi(z|x) = \mathcal{N}(z | \mu_\phi(x), \mathrm{diag}(\sigma_\phi(x))), \quad p_\theta(x|z) = \mathcal{N}(x | \mu_\theta(z), \mathrm{diag}(\sigma_\theta(z)))
\]

The reparametrization trick enables gradient-based optimization.

\subsection{Semi-Supervised Learning with VAEs}

Incorporating both labeled and unlabeled data:

Supervised loss (on labeled pairs):
\[
\mathcal{L}(x_i, y_i) = \mathbb{E}_{q(z|x_i, y_i)}[\log p(x_i, y_i|z) - \log q(z|x_i, y_i)]
\]

Unsupervised loss (on unlabeled):
\[
\mathcal{U}(x_j) = \mathbb{E}_{q(y,z|x_j)}[\log p(x_j, y|z) - \log q(z, y|x_j)]
\]

Combine with a classification term on labeled data to jointly optimize reconstruction and label prediction.

\subsection{Neural Networks as Gaussian Processes}

In the limit of infinite hidden-layer width with random initialization, a single-hidden-layer network:
\[
f(x) = a + w^T g(Vx + b)
\]

converges to a Gaussian process:
\[
f \sim \mathcal{GP}(0, k(x, x'))
\]

with kernel determined by the architecture:
\[
k(x, x') = \sigma_a^2 + \sigma_w^2 \mathbb{E}_{b,v}[g(v^T x + b) g(v^T x' + b)]
\]

\subsection{Information Theory Basics}

Mutual information measures dependence between random variables:
\[
I(X; Y) = D_{\mathrm{KL}}(p(x,y) \| p(x)p(y)) = H(X) - H(X|Y)
\]

For Gaussians, entropy is:
\[
H(X) = \frac{1}{2}\log \mathrm{det}(2\pi e \Sigma)
\]

\newpage

\section{Part VI: Advanced Applications and Inference Methods}

\subsection{Active Learning}

In active learning, we strategically select which unlabeled examples to label, maximizing model improvement per labeling cost:
\begin{itemize}
    \item Labeled set: $\{(x_i, y_i)\}_{i=1}^n$
    \item Unlabeled pool: $\{x_j\}_{j=1}^m$ with $m \gg n$
\end{itemize}

Simple strategies query points where model predictions are least confident (high entropy). More sophisticated approaches use Bayesian methods to measure uncertainty over the parameters.

\subsection{Bayesian Active Learning by Disagreement (BALD)}

BALD selects points that maximize parameter uncertainty by leveraging ensemble disagreement:
\[
x^* = \operatorname*{argmax}_x I(y; \theta | D, x) = H(y|D,x) - \mathbb{E}_{\theta \sim p(\theta|D)}[H(y|x,\theta)]
\]

The first term (average uncertainty) captures overall model uncertainty. The second term (expected conditional entropy) captures uncertainty that the model would retain after learning the label. The difference isolates epistemic (parameter) uncertainty.

\subsection{Bayesian Optimization}

For expensive black-box functions $f(x)$, Bayesian optimization enables efficient sequential evaluation:
\begin{enumerate}
    \item Fit a probabilistic surrogate model $p(f|D)$ (typically Gaussian Process)
    \item Select next evaluation point via acquisition function that balances exploration and exploitation
    \item Evaluate $f$ at selected point, update surrogate
    \item Repeat
\end{enumerate}

\subsection{Acquisition Functions}

Common acquisition functions trade off exploration and exploitation:

\textbf{Expected Improvement (EI):} Favors points likely to exceed current best:
\[
\alpha_{\mathrm{EI}}(x) = \mathbb{E}_f[\max(0, f(x) - f_{\text{best}})]
\]

\textbf{Upper Confidence Bound (UCB):} Balances mean and uncertainty:
\[
\alpha_{\mathrm{UCB}}(x) = \mu(x) + \beta \sqrt{v(x)}
\]

\textbf{Entropy Search:} Favors points that most reduce uncertainty about the optimum:
\[
\alpha_{\mathrm{ES}}(x) = H(x^*|D) - \mathbb{E}_{y|x}[H(x^* | D \cup \{(x,y)\})]
\]

\subsection{Langevin Dynamics for Posterior Sampling}

Rather than approximation, we can sample directly from the posterior using Langevin dynamics:
\[
\theta_{t+1} = \theta_t - \frac{\eta}{2} \nabla \log p(\theta_t|D) + \sqrt{\eta} \epsilon_t, \quad \epsilon_t \sim \mathcal{N}(0, I)
\]

With appropriate step size decay, this converges to samples from the posterior. The gradient term drives movement towards high-probability regions while noise enables exploration and prevents convergence to local modes.

\subsection{Stochastic Gradient Langevin Dynamics (SGLD)}

SGLD scales Langevin dynamics to large datasets using minibatches:
\[
\theta_{t+1} = \theta_t - \frac{\eta_t}{2} \nabla \tilde{L}(\theta_t) + \sqrt{\eta_t} \epsilon_t
\]

where $\tilde{L}(\theta) = \frac{N}{M} \sum_{i \in \text{batch}} \ell_i(\theta)$ rescales the minibatch loss. The noise from two sources---the minibatch gradient noise and the added Gaussian-drive efficient posterior exploration. As step sizes decay appropriately ($\eta_t \to 0$), SGLD asymptotically samples from the posterior.

\subsection{Wake-Sleep Algorithm}

The wake-sleep algorithm alternates between two phases for training deep generative models:

\textbf{Wake phase:} Update decoder using real data and learned encoder:
\[
\theta \gets \theta - \eta \nabla_\theta \mathbb{E}_{q_\phi(z|x)}[-\log p_\theta(x|z)]
\]

Improves model likelihood on real data.

\textbf{Sleep phase:} Update encoder using generated data from the decoder:
\[
\phi \gets \phi - \eta \nabla_\phi \mathbb{E}_{p_\theta(x,z)}[-\log q_\phi(z|x)]
\]

Improves encoder agreement with the model's implicit posterior.

This alternating approach avoids the KL divergence asymmetry issues present in standard VAE training and enables exploration of multimodal posterior structures. The algorithm naturally handles scenarios where a simple factorized Gaussian encoder would be insufficient.

\end{document}
